\emph{(i).} The pair $(M_{u}, M_{d})$ is obtained from $(Y, \tilde{Y})$ by the $2 \times 2$ array of scalars $\bigl[\begin{smallmatrix}\alpha & \bar{\beta} \\ \beta & \bar{\alpha}\end{smallmatrix}\bigr]$ (\claimref{synthesis.rung-field-coupling-to-its-conjugate-is-the-second-family-matrix}; the same formulas in the chart of Fact \claimref{synthesis.imported-left-right-bidoublet-two-yukawa-couplings}), whose determinant is $|\alpha|^{2} - |\beta|^{2} = 2m(c)$. For $m(c) \neq 0$ its inverse gives the stated formulas: $\bar{\alpha}M_{u} - \bar{\beta}M_{d} = (|\alpha|^{2} - |\beta|^{2})Y$ and $\alpha M_{d} - \beta M_{u} = (|\alpha|^{2} - |\beta|^{2})\tilde{Y}$. The map and its inverse are linear combinations with scalar coefficients, so they preserve symmetry of the matrices. Given $(M_{u}, M_{d})$ and a unit $z$, the pair $(zM_{u}, M_{d})$ has $H_{u}' = z\bar{z}H_{u} = H_{u}$ and $H_{d}' = H_{d}$; the masses are the square roots of the eigenvalues of $H_{q}$, the mixing matrix is built from their eigenvectors, and $D$ is a function of $H_{u}$, $H_{d}$ (Fact \claimref{synthesis.imported-jarlskog-commutator-identity}); $\det(zM_{u}) = z^{N}\det M_{u}$. As $z$ runs over the units $z^{N}$ runs over all of them. If $|\alpha| = |\beta|$ then $\bar{\alpha} = |\beta|^{2}/\alpha$, so $M_{d} = (\beta/\alpha)(\alpha Y + \bar{\beta}\tilde{Y}) = (\beta/\alpha)M_{u}$, $H_{d} = H_{u}$ and $D = 0$. The check verifies the inverse on a real and a non-real frame with symmetric samples, the turn by $z = \tfrac15(3 + 4I_{4})$ with $z^{3} = \tfrac{1}{125}(-117 + 44I_{4})$, and the balanced frame $\alpha = 3 + 4I_{4}$, $\beta = 5I_{4}$.

\emph{(ii).} Under $\alpha \mapsto z\alpha$, $\beta \mapsto \bar{z}\beta$: $M_{u} \mapsto zM_{u}$, $M_{d} \mapsto \bar{z}M_{d}$ and $\Theta \mapsto (z\bar{z})^{N}\Theta = \Theta$. Under the common phase the two terms of each $M_{q}$ turn oppositely, and $\Theta$ changes for generic matrices (checked on the sample). Under the exchange, $M_{u} \mapsto \alpha\tilde{Y}^{*} + \bar{\beta}Y^{*} = (\bar{\alpha}\tilde{Y} + \beta Y)^{*} = M_{d}^{*}$ and likewise $M_{d} \mapsto M_{u}^{*}$.

\emph{(iii).} (a) Real matrices have a real determinant, and real $M_{q}$ give real symmetric $H_{q}$, for which $D = 0$ (\claimref{synthesis.flag-conjugate-antilinear-on-one-chirality-and-conjugation-reverses-jarlskog}(iii)); changing the sign of $\tilde{Y}$ in the sample changes the sign of $\Theta$. (b) By (ii) with $\tilde{Y} = Y^{*}$ the pair is fixed by the exchange, so $M_{d} = M_{u}^{*}$, $\det M_{d} = (\det M_{u})^{*}$ and $\Theta = |\det M_{u}|^{2}$; $H_{d} = H_{u}^{*} = H_{u}^{\mathsf{T}}$ has the characteristic polynomial of $H_{u}$. With $H_{u} = A + I_{4}B$, $A$ real symmetric and $B$ real antisymmetric, $[H_{u}, H_{u}^{*}] = -2I_{4}[A, B]$ and $D = 8I_{4}\det[A, B]$ for $N = 3$, not zero in general; the sample has $D = 10909000000\,I_{4}$. (c) A real combination of Hermitian matrices is Hermitian and has a real determinant; the sample has $D \neq 0$ and different characteristic polynomials of $H_{u}$, $H_{d}$, and with a non-real frame $\Theta$ is not real. If $M = M^{\dagger} = M^{\mathsf{T}}$ then $M^{*} = M$.

\emph{(iv).} The condition $P_{12}XP_{12} = -X$ exchanges the indices $1$ and $2$ and changes the sign, which gives the four entries $p, q, r, s$ and $X_{33} = 0$. With $A = E + X$, $\det(\J + \delta A) = \delta^{3}\det A + \delta^{2}\sum_{ab}\mathrm{cof}(A)_{ab}$, the terms of order $\delta$ and $1$ vanishing because $\J$ has rank one. Expanding, $\det A = q^{2} - p^{2}$ and $\sum_{ab}\mathrm{cof}(A)_{ab} = q^{2} - p^{2} + 2r(p - q) + 2s(p + q) - 4rs = -(p - q - 2s)(p + q - 2r)$. In the basis $\hat{b}_{e} = (1, 1, -2)/\sqrt6$, $\hat{b}_{o} = (1, -1, 0)/\sqrt2$ of the cited sum rule, $x = \hat{b}_{e}^{\mathsf{T}}X\hat{b}_{o} = (p - q - 2s)/\sqrt3$ and $y = \hat{b}_{o}^{\mathsf{T}}X\hat{b}_{e} = (p + q - 2r)/\sqrt3$, so the sum of cofactors is $-3xy$: three times the determinant $-xy$ of the light-plane block (\claimref{synthesis.solar-angle-light-plane-sum-rule}(i)). The check verifies the formula on three samples with entries in $\mathrm{span}\{1, I_{4}\}$. For a symmetric one-type stage, $q = 0$, $s = r$ and the entries are $t$ times real numbers: $\det = -\delta^{2}t^{2}\bigl((p - 2r)^{2} + \delta p^{2}\bigr) = -\delta^{2}t^{2}R$, and $\det M_{q} = c_{q}^{3}$ times it. The product over the two sectors gives $\Theta$, the two signs cancelling. $R_{q}$ is real for real $\delta_{q}$, positive for $\delta_{q} > 0$, and for $\delta_{q} < 0$ of either sign: $p = 2$, $r = -1$, $\delta = -5$ gives $R = 16 - 20 = -4$, and $\delta = -4$ gives $R = 0$, a vanishing determinant. For real $c_{q}$ and nonzero $R_{q}$, $(c_{u}c_{d})^{3}(\delta_{u}\delta_{d})^{2}R_{u}R_{d}$ is a real number of sign $\sigma$. The off-diagonal entry of $B_{q}B_{q}^{\dagger}$ is $\tfrac23\bar{y}_{q} = \tfrac23\bar{t}_{q}(p_{q} - 2r_{q})/\sqrt3$, so the relative phase $\phi$ of \claimref{synthesis.cabibbo-from-one-type-quark-stages-selects-near-quarter-turn-phase} is $\pm(\arg t_{u} - \arg t_{d})$ up to a half turn from the signs of the real factors; $e^{2I_{4}\phi}$ is unaffected by the half turn. If $t_{d}$ is real, $(t_{u}t_{d})^{2}$ is a positive number times $e^{\pm2I_{4}\phi}$, which is real exactly when $\sin2\phi = 0$; $\cos2\phi = 2\cos^{2}\phi - 1 < -\tfrac12$ when $\cos^{2}\phi < \tfrac14$. At the quarter turn $(t_{u}t_{d})^{2} < 0$, so $\Theta$ has the sign $-\sigma$, which is $-\operatorname{sgn}(c_{u}c_{d})$ when both $\delta_{q} > 0$. With $t_{u} = I_{4}\tau$: $M_{u} = c_{u}(\J + \delta_{u}E) + I_{4}c_{u}\delta_{u}\tau S_{u}$, the first part even and the second odd under $P_{12}$, both real, so $P_{12}M_{u}^{*}P_{12} = M_{u}$; and $(M_{u}P_{12})^{\dagger} = P_{12}M_{u}^{*} = M_{u}P_{12}$, using that $M_{u}$ is symmetric: $M_{h} = M_{u}P_{12}$ is Hermitian and $M_{u} = M_{h}P_{12}$, the twin relation of \claimref{synthesis.one-type-second-stages-are-twins}, with $\det M_{u} = -\det M_{h}$. If one antiunitary conjugation $M \mapsto PM^{*}P$ with a real orthogonal $P$ fixes both matrices, then $H_{q} = PH_{q}^{*}P$ for both and $D = \det[H_{u}^{*}, H_{d}^{*}] = -D$ by clause (iii) of the result on the flag conjugate, so $D = 0$: with both coefficients real $P = 1$, with both at a quarter turn $P = P_{12}$. The check verifies on a sample with $t_{u} = 2I_{4}$, $t_{d} = 3$, $\delta_{u} = \tfrac15$, $\delta_{d} = \tfrac13$: $\Theta$ real and negative for $c_{u}c_{d} > 0$, positive for $c_{u}c_{d} < 0$; with $\delta_{u} = -5$, where $R_{u} = -4$, $\Theta = \tfrac{83417600}{3} > 0$ for $c_{u}c_{d} > 0$, the sign $-\sigma$ and not $-\operatorname{sgn}(c_{u}c_{d})$; with $\delta_{u} = -4$, $\Theta = 0$; $D \neq 0$, different characteristic polynomials, the two conjugation properties and the Hermitian twin; that an eighth turn against a real coefficient gives a non-real $\Theta$ and opposite eighth turns a positive one; and that $\Theta/(t_{u}t_{d})^{2}$ is real with the sign of $c_{u}c_{d}$ on two generic pairs.

\emph{Reading.} The action clause in its adjoint reading takes each family term to its adjoint; the action is the sum of each term and its adjoint for every $(Y, \tilde{Y})$, so the clause holds identically. The other statements of the Reading combine (i), (iii), (iv) with the stored exclusions named there. That a physical angle of a half turn conserves the conjugation: the conjugation reverses the angle, and a half turn equals its own reverse modulo a whole turn. The neutron dipole bound is quoted in the store through a conversion linear in the angle, valid near zero; it gives no bound near a half turn.
